\section{Trabalhos Relacionados}

\subsection{Canais e Estados de Substituição em Caches}

A pesquisa \textit{“Last-Level Cache Side-Channel Attacks are Practical”}\cite{7163050}, demonstrou o quão viável é explorar caches compartilhados, como a LLC - Last-Level Cache, em ambientes multi threaded e de nuvem para extrair chaves criptográficas. No entanto, enquanto os ataques clássicos focam no tempo do acesso à memória, o artigo em questão refina essa abordagem ao demonstrar que a própria política de substituição LRU (Least Recently Used) retém estado e é capaz do vazamento de informação. Como limitação, o ataque depende do uso de \textit{large pages} pelo atacante para montar os conjuntos de desvio (\textit{eviction sets}), e os próprios autores reconhecem que desabilitar esse recurso na VMM seria a defesa mais simples. Além disso, a arquitetura de LLC dividida em \textit{slices} obriga os autores a montar esses conjuntos por tentativa e erro. O ataque também exige etapa manual de análise e os autores admitem não ter testado o método em ambiente ruidoso ou em nuvem real.

Ademais, a análise dos componentes e estados internos do processador se conecta com \textit{“Are Coherence Protocol States Vulnerable to Information Leakage?”}\cite{8327007}, porque ambas as pesquisas compartilham a mesma ideia de que estruturas de otimização de hardware são capazes de vazarem dados através de suas alterações de estado, não se limitando ao tempo de transferência de dados. Contudo, os autores mostram que, acima de 500~Kbps, a maioria das configurações testadas perde precisão significativamente. O ataque também exige cooperação entre dois processos maliciosos, trojan e spy, o que caracteriza um canal encoberto e restringe seu uso a cenários de vazamento interno. Além disso, a criação de memória compartilhada via KSM leva dezenas de segundos, e os testes foram feitos em um único modelo de processador.

Em \textit{“Leaking Information Through Cache LRU States”}\cite{9065598} introduz-se a exploração da LRU, revelando falhas em arquiteturas que eram projetadas para isolar acessos. Porém, como qualquer acesso ao conjunto de cache altera o estado LRU, o receptor só pode medir o conjunto uma vez por rodada, obtendo no máximo um bit por medição, o que é uma limitação estrutural do canal. A política PLRU agrava isso, gerando erros de leitura. A técnica de medição por \textit{pointer chasing} também exige calibrar o tamanho da lista ligada. Por fim, a taxa de transmissão varia bastante conforme o cenário (600~Kbps em \textit{hyper-threading} contra poucos bps em compartilhamento por tempo).

Esses trabalhos adotam metodologias bem distintas entre si. \cite{7163050} avalia ataques na LLC compartilhada, implicando em modelos de ameaça onde atacar a LLC permite ataques entre VMs distintas, enquanto \cite{9065598} concentra-se no estado LRU do cache L1 privado, o que normalmente exige coexecução no mesmo núcleo. Quanto à técnica de medição, \cite{7163050} usam \textit{Prime+Probe}, \cite{9065598} usam uma variação de \textit{pointer chasing} e \cite{8327007} medem diretamente a latência de acesso a blocos em diferentes estados de coerência.

\subsection{Mitigações e Lacunas da Literatura}

No escopo das defesas, \textit{“CATalyst: Defeating last-level cache side channel attacks in cloud computing”}\cite{7446082} propõe o uso da tecnologia Cache Allocation Technology (CAT) do hardware para particionar a LLC e neutralizar ataques. O artigo dialoga com essa solução ao tratar a eficácia de mecanismos de isolamento, haja visto que as defesas puramente baseadas em software ou particionamento parcial podem ser contornadas. A limitação é que a solução depende da tecnologia Intel CAT, herdando o limite de apenas 4 \textit{Classes of Service}, insuficiente para isolar muitas VMs. O particionamento fino por página também tem limites próprios, como poucas páginas seguras por VM e perda de desempenho. Além disso, o CATalyst neutraliza ataques por contenção de linha, mas não os canais baseados em metadados de coerência ou de substituição LRU explorados em~\cite{8327007} e~\cite{9065598}.

Os ambientes experimentais das metodologias também revelam lacunas: trabalhos como \cite{7163050} e \cite{7446082} testam em máquinas virtualizadas simulando nuvem, enquanto \cite{8327007} e \cite{9065598} testam em processos nativos sem virtualização. 

Desta forma, embora Xiong et al.~\cite{9354935} tenham demonstrado a viabilidade de canais encobertos baseados no estado de substituição LRU, a avaliação de contramedidas na literatura carece de uma base comparativa quantitativa e unificada. A pesquisa original foca sua avaliação empírica via simulação (gem5) apenas nas defesas de \textit{Partition-Locked (PL) cache} e \textit{Random Fill (RF) cache}. Outras abordagens estado-da-arte, como DAWG, ScatterCache, CEASER e InvisiSpec, tiveram sua resiliência analisada de forma estritamente qualitativa. Há uma lacuna significativa na literatura quanto à medição empírica da banda residual do canal encoberto frente a essas defesas.

Além disso, estudos anteriores assumem que canais baseados em LRU em níveis mais baixos da hierarquia de memória (L2/LLC) seriam facilmente mitigáveis por técnicas preexistentes em L1. Tal suposição é feita sem uma validação experimental rigorosa e desconsidera o comportamento de políticas de substituição modernas comumente empregadas na LLC, como RRIP e suas variantes.

Para preencher essas lacunas dentro de um escopo viável, este trabalho propõe uma avaliação experimental focada, utilizando o simulador gem5 em modo \textit{syscall-emulation} (SE). Nossa metodologia adapta os protocolos de canais com e sem memória compartilhada (Algoritmos 1 e 2 de~\cite{9354935}) para medir métricas de segurança em conjunto com métricas de desempenho. 

Em suma, este artigo apresenta as seguintes contribuições principais:
\begin{itemize}
    \item \textbf{Validação Empírica na LLC:} A avaliação da capacidade do canal de substituição contra políticas modernas na LLC (como a família RRIP, nativa ou adaptada no ecossistema gem5), testando a suposição teórica da literatura de que camadas mais baixas da hierarquia seriam inerentemente mais fáceis de proteger;
    \item \textbf{Análise de Desempenho Integrada:} Uma avaliação do \textit{trade-off} entre segurança e desempenho, medindo o impacto em IPC sob cargas de trabalho abertas e controladas (ex.: PARSEC, MiBench ou microbenchmarks sintéticos);
    \item \textbf{Extensão Comparativa (Opcional):} Caso a curva de aprendizado da ferramenta permita, a implementação e medição quantitativa de uma defesa por particionamento dinâmico do estado de substituição (ex.: adaptação do conceito do DAWG) para fins de comparação direta com os dados originais.
\end{itemize}